\section{Spatial parity of traces}\label{v11:sec:parity-limits}

Reflection sectors have classical determinant analogues in scalar
Toeplitz-plus-Hankel theory \cite[Sections~2--3]{BasorEhrhardt2002}.
The following argument treats the operator-valued symbol and its seam
directly, retaining the physical within-cell reflection.

\begin{lemma}[A reflected Toeplitz limit]
\label{ext:lem:reflected-toeplitz}
Let $r$ be a self-adjoint involution on a separable Hilbert space
$\mathscr H$, and let $a\in C^3([0,1];\mathfrak S_1)$ be self-adjoint
with $ra(1-\theta)r=a(\theta)$.  Put
\[
 A=\tfrac12(a(0)+a(1)),\qquad B=\tfrac12(a(0)-a(1)),
 \qquad R_n=F_n\otimes r,
\]
where $F_n$ reverses the $n$ cells and $M=\sup_\theta\|a(\theta)\|$.
For every $f(0)=0$ holomorphic near $[-M,M]$, along even $n$,
\begin{equation}\label{ext:eq:reflected-toeplitz}
 \lim_{n\to\infty}\operatorname{tr}(R_nf(T_n(a)))
 =\tfrac12\operatorname{tr}_{\mathscr H}
       \bigl(r[f(A)-f(a(1/2))]\bigr).
\end{equation}
\end{lemma}

\begin{proof}
All finite sections are trace class: their trace norms are at most
$n\int_0^1\|a(\theta)\|_1\,d\theta$, by the rank-one cell-vector
representation of the Fourier coefficients.  Set
$h(\theta)=1-2\theta$, $b=A+hB$, $c=a-b$, and $a_t=b+tc$.
Then $c(0)=c(1)=0$, $R_n$ commutes with $T_n(a_t)$, and
$\|a_t(\theta)\|\le M$: both $b$ and $a_t$ are convex combinations
of symbols with norm at most $M$.

Put $H_n=T_n(h)$. Since $F_nH_nF_n=-H_n$ and
$T_n(b)=I_n\otimes A+H_n\otimes B$, reflection exchanges the
nonzero $H_n$ eigenspaces in opposite pairs, contributing zero to
$\tr(R_nf(T_n(b)))$. The kernel contributes
$\tr(F_n|_{\ker H_n})\tr(rf(A))=0$: for even $n$,
$\tr F_n=0$ and its nonzero pairs have zero trace. Hence
\begin{equation}\label{ext:eq:linear-reflection-zero}
 \tr(R_nf(T_n(b)))=0.
\end{equation}

If $d\in C^3([0,1];\mathfrak S_1)$ and $d(0)=d(1)$,
Lemma~\ref{rev5:lem:seam} makes both one-sided Hankel blocks
$\Pi_+\mathcal L(d)\Pi_-$ and $\Pi_-\mathcal L(d)\Pi_+$ trace class:
the principal matrix is $((j+l+1)^{-2})_{j,l\ge0}$ and the remainder
has summable weighted Fourier trace norms. For a fixed bounded
self-adjoint symbol $a_0$ and $k\ge0$, this gives
\begin{equation}\label{ext:eq:reflected-boundary-zero}
 \tr\bigl(R_n[T_n(d)T_n(a_0)-T_n(da_0)]T_n(a_0)^k\bigr)\to0.
\end{equation}
Indeed split $-P_n\mathcal L(d)(I-P_n)\mathcal L(a_0)P_n$ into
$j<0$ and $j\ge n$ boundary terms. At the left, the Hankel factor
converges in trace norm and all remaining factors and adjoints converge
strongly. Its trace-class limit pairs to zero with the weakly null
zero-extended $R_n$ (check finite cell supports). Conjugating by $F_n$
gives the same argument at the right, since $F_nR_nF_n=R_n$.
Finite-rank approximation makes these bounds uniform on compact
$C^3$ symbol families.

Choose a fixed contour $\Gamma$ in the holomorphy domain enclosing
$[-M,M]$ at positive distance.  For $\zeta\in\Gamma$, put
$d_{t,\zeta}=c(\zeta-a_t)^{-1}$.  These symbols are $C^3$ in
trace norm, vanish at both endpoints, and have uniform $C^3$ bounds
on $[0,1]\times\Gamma$.  With $Q_n=I-P_n$, exactly
\[
 T_n(c)(\zeta-T_n(a_t))^{-1}-T_n(d_{t,\zeta})
 =-P_n\mathcal L(d_{t,\zeta})Q_n\mathcal L(a_t)P_n
           (\zeta-T_n(a_t))^{-1}.
\]
Put $\delta=\operatorname{dist}(\Gamma,[-M,M])$.
The resolvents have norm at most $\delta^{-1}$; strong convergence
of the self-adjoint finite sections and the resolvent identity give
strong convergence of resolvents and adjoints at both ends. Finite
nets make this uniform on fixed vectors for $(t,\zeta)\in[0,1]\times\Gamma$.
The symbols $d_{t,\zeta}$ form a compact $C^3(\mathfrak S_1)$ family,
so their two Hankel trace norms have a common bound $C_H$.
The displayed error thus has trace norm at most $MC_H/\delta$.
Its boundary limits are trace-norm uniform by the preceding localization;
finite-rank approximation and weakly null reflection give uniform
vanishing of its reflected trace. Cauchy calculus therefore gives
\[
 \operatorname{tr}(R_nT_n(c)f'(T_n(a_t)))
 -\operatorname{tr}(R_nT_n(cf'(a_t)))\longrightarrow0.
\]
Differentiating the Cauchy integral and cycling the trace gives
\[
 \frac d{dt}\operatorname{tr}(R_nf(T_n(a_t)))
 =\operatorname{tr}(R_nT_n(c)f'(T_n(a_t))).
\]
For $e_t=cf'(a_t)$ the zero-endpoint Fourier sums are uniformly
absolutely summable in trace norm, say by $C_e$. Thus the derivative
integrand is dominated, independently of $n,t$, by
\[
 C_e+\frac{\operatorname{length}(\Gamma)}{2\pi\delta}
                    \sup_\Gamma|f'|\,MC_H.
\]
This justifies the contour limit and dominated convergence in $t$.
For even $n$, absolute trace-norm Fourier convergence gives
\[
 \tr(R_nT_n(e_t))=\sum_{|k|<n,\ k\text{ odd}}\tr(r\widehat e_t(k))
 \longrightarrow\tfrac12\tr(r[e_t(0)-e_t(1/2)]).
\]
Reflection symmetry gives $[r,a_t(1/2)]=0$, so cyclicity identifies
$\tr(rc(1/2)f'(a_t(1/2)))$ with $d\tr(rf(a_t(1/2)))/dt$.
Since $e_t(0)=0$, integration in $t$ gives
$\tfrac12\tr(r[f(A)-f(a(1/2))])$. This proves
\eqref{ext:eq:reflected-toeplitz} for this $f$.  The exact linear
reference cancellation works unchanged by its spectral decomposition.


\end{proof}

\begin{theorem}[Spatial-parity limit]
\label{ext:thm:polynomial-parity}\label{ext:thm:holomorphic-parity}
Let $Rf(x)=f(-x)$ on $L^2[-1,1]$ and $g=1/\sqrt2$.
For every $f(0)=0$ holomorphic near $[-1,1]$,
\begin{equation}\label{ext:eq:polynomial-parity}
 \operatorname{tr}(Rf(\mathcal K_\tau))
 =\tfrac12[f(g)+f(-g)+f(1)-f(-1)]+o_f(1)
 \qquad(\tau\to\infty).
\end{equation}
In particular, for every fixed integer $m\ge1$,
\[
 \operatorname{tr}(R\mathcal K_\tau^{2m})\longrightarrow2^{-m},
 \qquad
 \operatorname{tr}(R\mathcal K_\tau^{2m-1})\longrightarrow1.
\]
\end{theorem}

\begin{proof}
At integer $\tau=q$, dilation and integer translation identify the
operator with the $2q$-cell section of its exact symbol
$\sigma(\theta)=W_\theta^*X_{1,0}W_\theta$.
Spatial reflection becomes $F_{2q}\otimes r$, where
$(rf)(t)=f(1-t)$.  This symbol is smooth in trace norm and satisfies
$r\sigma(1-\theta)r=\sigma(\theta)$.
Its endpoint average
$A=(X_{1,0}+X_{0,-1})/2$ has nonzero eigenvalues $\pm g$,
both in the $r$-even space
$\operatorname{span}\{e_0,(e_1+e_{-1})/\sqrt2\}$.
At $\theta=1/2$, the $+1$ eigenvector of $\sigma$ is proportional
to $\cos(\pi t)$ and is $r$-odd, while the $-1$ eigenvector is
proportional to $\sin(\pi t)$ and is $r$-even.  Hence
Lemma~\ref{ext:lem:reflected-toeplitz} gives precisely
\eqref{ext:eq:polynomial-parity} along integer $\tau$.

For $\tau=n+r$, $r\in[0,1]$, set $m=2n$,
$A_m=P_{(-r,m+r)}\mathcal G P_{(-r,m+r)}$ and
$B_m=P_{(0,m)}\mathcal G P_{(0,m)}$; both commute with
$\mathscr R_m u(X)=u(m-X)$. The already proved strip localization
\eqref{v11:eq:strip-localization} gives
\[
 A_m-B_m=F_{-r}^L+\mathsf S^mF_r^R\mathsf S^{-m}+o_{\mathfrak S_1}(1).
\]
For a contour $\Gamma$ around $[-1,1]$ in the holomorphy domain, insert
this identity into
\[
 f(A_m)-f(B_m)=\frac1{2\pi i}\int_\Gamma f(\zeta)
 (\zeta-A_m)^{-1}(A_m-B_m)(\zeta-B_m)^{-1}\,d\zeta.
\]
Each boundary product converges in trace norm after fixing that
boundary.  The reflections $\mathscr R_m=\mathsf S^mR$ are weakly null
at the left, and $\mathsf S^{-m}\mathscr R_m\mathsf S^m=\mathsf S^{-m}R$ are weakly null
at the right.  Consequently
$\operatorname{tr}(\mathscr R_m[f(A_m)-f(B_m)])\to0$ for fixed $r$.

Cauchy calculus bounds $\|f(A)-f(B)\|_1$ by a fixed multiple of
$\|A-B\|_1$, while
$\|A_m(r)-A_m(s)\|_1\le4b(|r-s|)$ independently of $m$.
Thus finite nets give uniform convergence on $r\in[0,1]$.
Dilation and integer translation recover the reflected traces at $n+r$
and $n$, proving the full limit.
\end{proof}

\begin{remark}[Quantitative moment transfer]
Writing $G_L=P_{(-L,L)}\mathcal G P_{(-L,L)}$, cyclic differentiation
of the reflected multiple-integral trace gives
$d\tr(RG_L^m)/dL=2mG_L^m(L,-L)$: all $m$ boundary terms agree,
and the two endpoints agree by reflection. With $w(s)=(1+|s|)^{-1}$,
$|G(X,Y)|\le4w(X-Y)$ and
$\int_{-L}^Lw(X-Z)w(Z-Y)dZ\le8\log(1+2L)w(X-Y)$.
Indeed one of the two distances is at least $|X-Y|/2$, while the
other factor has integral at most $2\log(1+2L)$. Induction gives
\[
 |G_L^m(L,-L)|\le\frac{4^m[8\log(1+2L)]^{m-1}}{1+2L}.
\]
Thus replacing $L$ by $\lfloor L\rfloor$ changes its reflected
$m$th moment by $O_m(\log^{m-1}(2+L)/(1+L))$.
\end{remark}

\begin{corollary}[Parity trace balance and determinants]
\label{ext:cor:polynomial-parity-sectors}
\label{ext:cor:parity-determinants}
For the spatial even and odd restrictions $\mathcal K_{\tau,\pm}$ and
$f(0)=0$ holomorphic near $[-1,1]$,
\[
 \operatorname{tr}f(\mathcal K_{\tau,\pm})
 =\tfrac12\operatorname{tr}f(\mathcal K_\tau)
 \ \pm\tfrac14[f(g)+f(-g)+f(1)-f(-1)]+o_f(1).
\]
Thus the volume and logarithmic coefficients divide equally.
For each real $w>0$,
\begin{equation}\label{ext:eq:parity-squared-determinant}
 \frac{\det(I+w\mathcal K_{\tau,+}^{2})}
      {\det(I+w\mathcal K_{\tau,-}^{2})}
 \longrightarrow1+\frac w2.
\end{equation}
For $|z|<1$, locally uniformly, the analytic logarithms normalized at zero
satisfy
\begin{equation}\label{ext:eq:parity-linear-determinant}
 \log\frac{\det(I-z\mathcal K_{\tau,+})}
               {\det(I-z\mathcal K_{\tau,-})}
 \longrightarrow\tfrac12\log(1-z^2/2)
             +\tfrac12\log\frac{1-z}{1+z}.
\end{equation}
\end{corollary}
\begin{proof}
The commuting projections $(I\pm R)/2$ and
Theorem~\ref{ext:thm:holomorphic-parity} give the first identity.
Apply it to $\log(1+wx^2)$, holomorphic near $[-1,1]$, and
$\log(1-zx)$ to obtain the determinant formulas.
For $|z|\le\rho<1$ choose a common contour $\Gamma$ around $[-1,1]$
with $\rho\sup_\Gamma|\zeta|<1$. The normalized logarithms and their
derivatives form a compact family there; the proof's Hankel, Fourier
and strip bounds and finite nets give local uniformity.
\end{proof}

\subsection{Sharp parity counts inside the central gaps}
\label{boundary:sec:sharp-parity}

\begin{theorem}[Exact eventual parity counts in the central gaps]
\label{boundary:thm:sharp-parity}
Fix $r\in\mathbb R$ and $0<t<g$. If $t\notin\operatorname{spec}B_{-r}$,
then, for all sufficiently large integers $n$,
\begin{equation}\label{boundary:eq:positive-sharp-parity}
 \#\{\lambda_j(\mathcal K_{n+r,+})>t\}
 -\#\{\lambda_j(\mathcal K_{n+r,-})>t\}=1.
\end{equation}
If instead $-t\notin\operatorname{spec}B_{-r}$, then eventually
\begin{equation}\label{boundary:eq:negative-sharp-parity}
 \#\{\lambda_j(\mathcal K_{n+r,+})<-t\}
 -\#\{\lambda_j(\mathcal K_{n+r,-})<-t\}=0.
\end{equation}
Here the subscripts $+$ and $-$ denote even and odd spatial parity.
Each assertion requires regularity only at its own signed threshold.
\end{theorem}

\begin{proof}
Put $m=2n$ and use physical coordinates $(-r,m+r)$. Write
$\rho f(t)=f(1-t)$ on the fiber and
$\mathcal R_m=F_m\otimes\rho$ on the full cell space, where the full
cell reflection sends $j$ to $m-1-j$. Retain the gauge of Lemma~\ref{rev5:lem:gauge}
$\mathscr U=\mathcal L(V)$ and set
\[
 \widetilde{\mathcal R}_m=\mathscr U^*\mathcal R_m\mathscr U,
 \quad Z(\theta)=\rho V(1-\theta)^*\rho V(\theta),\quad d=Z-I.
\]
Then
\begin{equation}\label{boundary:eq:reflection-gauge}
 \widetilde{\mathcal R}_m=\mathcal R_m\mathcal L(Z),\qquad
 \rho Z(1/2)=V(1/2)^*\rho V(1/2),
\end{equation}
and $d$ is smooth in trace norm with $d(0)=d(1)=0$.

Let $P_m^{A,\sigma}$ and $P_m^{C,\sigma}$ be the sharp spectral
projections of $\mathscr A_m(-r,r)$ and $\mathscr C_m$ above $t$
for $\sigma=+$, or below $-t$ for $\sigma=-$.
Choose a rectangle enclosing the appropriate spectral band, crossing
the real axis at the regular threshold and beyond $1$ or $-1$.
The gap theorem gives uniform contour resolvent bounds for large $m$
and both half-line limits. Their resolvents and adjoints converge
strongly at each end, including the real crossing by the resolvent
identity and common bound. Inserting
\eqref{v11:eq:combined-gap-localization} therefore localizes
$P_m^{A,\sigma}-P_m^{C,\sigma}$ in trace norm at both boundaries.
The gauge commutes with shifts, so both
$\widetilde{\mathcal R}_m$ and
$\mathsf S^{-m}\widetilde{\mathcal R}_m\mathsf S^m$ are weakly null.
Consequently
\begin{equation}\label{boundary:eq:sharp-boundaries-vanish}
 \operatorname{tr}\bigl(\widetilde{\mathcal R}_m
 (P_m^{A,\sigma}-P_m^{C,\sigma})\bigr)\longrightarrow0.
\end{equation}

The star itself has exactly zero reflected count in each sign for even
$m$. To see this, write it as
$\left(\begin{smallmatrix}0&B_m\\B_m^*&0\end{smallmatrix}\right)$,
where $B_mB_m^*\ge I/2$, and put
$W_m=(B_mB_m^*)^{-1/2}B_m$.
Its signed spectral subspaces are the graphs
$x\mapsto2^{-1/2}(x,\pm W_m^*x)$.
The star reflection symmetry intertwines these graphs with $F_m$ on
the central component. Since $\operatorname{tr}F_m=0$ and all nonzero
star eigenvalues have magnitude at least $g$,
$\operatorname{tr}(\mathcal R_mP_m^{C,\sigma})=0$ exactly.
Thus \eqref{boundary:eq:reflection-gauge} reduces the remaining trace to
\[
 \operatorname{tr}(\widetilde{\mathcal R}_mP_m^{C,\sigma})
 =\operatorname{tr}(\mathcal R_mT_m(d)P_m^{C,\sigma}).
\]

For $z$ on the contour put $d_z=d(z-a)^{-1}$. An exact identity is
\[
 T_m(d)(z-\mathscr C_m)^{-1}-T_m(d_z)
 =-P_m\mathcal L(d_z)(I-P_m)\mathcal L(a)P_m
                         (z-\mathscr C_m)^{-1}.
\]
The symbols $d_z$ have zero endpoints and trace-class one-sided Hankel
blocks, uniformly on the contour. The localization argument of
Lemma~\ref{ext:lem:reflected-toeplitz}, now using the reference resolvent,
makes the reflected trace of the right side tend to zero uniformly.
Contour integration and the odd Fourier coefficient sum therefore give
\[
 \operatorname{tr}(\mathcal R_mT_m(d)P_m^{C,\sigma})
 \longrightarrow-\tfrac12\operatorname{tr}_{\mathscr H}
 \bigl(\rho d(1/2)p_\sigma(a(1/2))\bigr),
 \qquad p_\sigma(a)=\tfrac12(a^2+\sigma a).
\]
Both signed star eigenvectors at $1/2$ are $\rho$-even. The physical
positive eigenvector is proportional to $\cos(\pi t)$ and is
$\rho$-odd; the negative one is proportional to $\sin(\pi t)$ and
is $\rho$-even. Equation~\eqref{boundary:eq:reflection-gauge} makes the
two fiber traces respectively $-2$ and $0$. Together with
\eqref{boundary:eq:sharp-boundaries-vanish}, the physical reflected
counts tend to $1$ and $0$. They are integers because physical
reflection commutes with the physical sharp spectral projections,
so both equalities hold eventually.
\end{proof}

In particular the gauge-invariant two-boundary positive count offset
is odd, and its negative counterpart is even. These congruences do not
fix an individual reference index or evaluate a boundary eigenvalue.

\begin{corollary}[Alternation of signed boundary magnitudes]
\label{interaction:cor:boundary-alternation}
There is a continuous strictly decreasing lift
$\mathfrak b:(0,g)\to\mathbb R$ such that, for every real $b$,
\[
 t\in\operatorname{spec}B_b
 \iff\mathfrak b(t)\in b+\mathbb Z,
 \qquad
 -t\in\operatorname{spec}B_b
 \iff\mathfrak b(t)\in b-\tfrac12+\mathbb Z.
\]
Thus the positive gap levels and the magnitudes of the negative gap
levels alternate at each fixed boundary phase. If $0<a<c<g$ and
$\pm a,\pm c\notin\operatorname{spec}B_b$, their counts in the
magnitude interval $(a,c)$ differ by at most one.
For the centered operators $\mathcal K_{n+r}$, these two signed counts
are eventually even and differ by at most two; in each parity sector
they differ by at most one.
\end{corollary}
\begin{proof}
The simple strictly decreasing branches of
Theorem~\ref{boundary:thm:one-cell-flow} have continuous strictly
decreasing local inverses. Unique phases make them agree modulo one;
the interval $(0,g)$ admits a continuous real lift. Subdivision of
compact subintervals makes it globally strictly decreasing.
Half-period sign reversal gives the second equivalence; differentiability
of the inverse is unnecessary.

The two lattices in the assertion alternate. More explicitly, at the
stated regular endpoints the two counts are
\[
 \lfloor\mathfrak b(a)-b\rfloor-\lfloor\mathfrak b(c)-b\rfloor,
 \quad
 \lfloor\mathfrak b(a)-b+\tfrac12\rfloor
       -\lfloor\mathfrak b(c)-b+\tfrac12\rfloor.
\]
Their difference has magnitude at most one. For $b=-r$, the centered
finite statement follows from Theorem~\ref{v11:thm:gap-limits} and
Corollary~\ref{boundary:cor:opposite-signs}: each boundary level
contributes exactly one eigenvalue of each parity, and there are
eventually no additional eigenvalues in either fixed compact interval.
\end{proof}

\begin{proposition}[Order of a stationary boundary crossing]
\label{interaction:prop:contact-order}
Let $\lambda_0\in(-g,0)\cup(0,g)$ be an eigenvalue of $B_{b_0}$,
and let $h_0$ be its normalized physical eigenfunction, continued to
its entire bandlimited representative. Let $d\ge0$ be its finite
vanishing order at $b_0$, with $d=0$ when $h_0(b_0)\ne0$.
The local eigenbranch satisfies
\[
 \lambda(b_0+s)=\lambda_0-
 \frac{\lambda_0|h_0^{(d)}(b_0)|^2}{(2d+1)(d!)^2}s^{2d+1}
 +O(s^{2d+2}).
\]
In particular, any stationary crossing has odd finite order.
\end{proposition}
\begin{proof}
Use the unshifted frequency pencil
$L(b)=Q_b-\lambda(b)J$ from
\eqref{boundary:eq:covariance-rank-one}. Choose an analytic normalized
kernel vector $u(b)$ with $\langle u,u'\rangle=0$, and set
\[
 z(b)=\langle v_b,u(b)\rangle,\quad
 q(b)=\langle u(b),Q_bu(b)\rangle>0,\quad
 j(b)=\langle u(b),Ju(b)\rangle=\frac{q(b)}{\lambda(b)}\ne0.
\]
Differentiating $L(b)u(b)=0$ gives
\[
 \lambda'=-\frac{|z|^2}{j},\qquad
 L(b)u'=v_bz+\lambda'Ju.
\]
The inverse of $L(b)$ on the orthogonal complement of its simple
kernel is locally uniformly bounded. Hence
$\|u'\|\le C(|z|+|z|^2)$.
The analytic function $z$ cannot vanish on an interval, since then
$\lambda$ would be constant. If its zero order at $b_0$ is $d_0$,
integration gives $u(b_0+s)-u(b_0)=O(s^{d_0+1})$.
Consequently $z(b_0+s)$ and
$\langle v_{b_0+s},u(b_0)\rangle$ have the same first nonzero Taylor
coefficient. The latter is $\sqrt{q(b_0)}h_0(b_0+s)$, so $d_0=d$ and
\[
 z(b_0+s)=\sqrt{q(b_0)}\frac{h_0^{(d)}(b_0)}{d!}s^d
             +O(s^{d+1}).
\]
Substitute this in the formula for $\lambda'$ and integrate once.
\end{proof}

\begin{theorem}[Relative signed counts at a regular gap threshold]
\label{interaction:thm:signed-count-imbalance}
Fix $r\in\mathbb R$ and $0<t<g$ with
$\pm t\notin\operatorname{spec}B_{-r}$. Write
\[
 N_+(t;n,r)=\#\{\lambda_j(\mathcal K_{n+r})>t\},\qquad
 N_-(t;n,r)=\#\{\lambda_j(\mathcal K_{n+r})<-t\}.
\]
For all sufficiently large integers $n$, putting $b=-r$, one has
\[
 N_+(t;n,r)-N_-(t;n,r)
 =2\bigl[I_+(t;b)-I_+(t;b+\tfrac12)\bigr]-1
 \in\{-1,1\}.
\]
The value is $1$ precisely when the unique positive crossing phase
lies in $(b,b+1/2)$ modulo one. The two physical count offsets relative
to $2n$ are therefore adjacent integers; their absolute values are
not determined here.
\end{theorem}
\begin{proof}
Write $i_L(t;l)=\iota_+(t;l)$ and
$i_R(t;h)=\iota_-(t;h)$ for the positive left and right indices in
\eqref{v11:eq:gap-indices}. The references cancel in their differences.
Since physical reflection commutes with $\mathcal G$ and exchanges the
right boundary at $h$ with the left boundary at $-h$, unitary invariance
of the trace-class projection difference gives
\[
 i_R(t;h_2)-i_R(t;h_1)
 =i_L(t;-h_2)-i_L(t;-h_1)
\]
at regular endpoints. This does not require a reflection-equivariant
gauge or an equality between individual reference indices.

Put $m=2n$. Half-translation changes the kernel's sign, so the negative
count on $(b,m-b)$ equals the positive count on
$(b+1/2,m-b+1/2)$. Apply
Corollary~\ref{v11:cor:gap-count-offset} to these two fixed endpoint
pairs, using the same $m$. Their reference counts cancel, and reflection
then gives
\begin{align*}
 N_+-N_-
 &=i_L(t;b)-i_L(t;b+\tfrac12)
     +i_R(t;-b)-i_R(t;-b+\tfrac12)\\
 &=2i_L(t;b)-i_L(t;b+\tfrac12)-i_L(t;b-\tfrac12)\\
 &=2\bigl[i_L(t;b)-i_L(t;b+\tfrac12)\bigr]-1.
\end{align*}
The period law and monotonicity in
Theorem~\ref{boundary:thm:one-cell-flow} make the bracket either zero
or one; its value is one exactly when the unique crossing occurs in
that half-period. Regularity at its endpoints follows from the two
signed regularity assumptions in the statement.
\end{proof}

\begin{corollary}[Relative ranks of a centered boundary pair]
\label{interaction:cor:relative-cluster-ranks}
Let $\alpha$ be a nonzero central-gap eigenvalue of $B_{-r}$.
Choose $0<a<|\alpha|<c<g$ so that $\alpha$ is the only boundary
eigenvalue of either sign with magnitude in $[a,c]$.
For all sufficiently large $n$, let $k$ be the number of opposite-sign
eigenvalues of $\mathcal K_{n+r}$ with magnitude greater than $c$.
The two eigenvalues converging to $\alpha$ occupy same-sign ranks
$k,k+1$ and combined absolute-value ranks $2k,2k+1$, counting from
largest magnitude and counting multiplicity.

If $\alpha>0$, then $k=2q$ and their even and odd positive-sector ranks
are respectively $q+1$ and $q$. If $\alpha<0$, then $k=2q+1$ and their
respective negative-sector ranks are both $q+1$.
\end{corollary}
\begin{proof}
Local gap convergence adds exactly two own-sign eigenvalues and no
opposite-sign eigenvalue between the thresholds. The latter count stays
$k$; the former increases by two. The signed-count theorem forces its
value above $c$ to be $k-1$, giving both rank pairs.

For a positive cluster, Theorem~\ref{boundary:thm:sharp-parity} makes
$k=2q$ and divides the $k-1$ positive levels above the cluster into
$q$ even and $q-1$ odd levels. For a negative cluster it makes $k=2q+1$
and divides the $k-1$ negative levels into $q$ of each parity.
Corollary~\ref{boundary:cor:opposite-signs} then adds one member in
each sector. The internal ordering of the two parities is not asserted.
\end{proof}



